%% Chapitre4.tex — Réalisation et résultats
\chapter{Réalisation et résultats}
\label{chapter:realisation}

\epigraphe{Les paroles s'envolent, le code reste.}{Adapté d'un proverbe latin}

%% ── Introduction ───────────────────────────────────────────────────────────
\section*{Introduction}
\addcontentsline{toc}{section}{Introduction}
\par Ce chapitre décrit la mise en œuvre concrète de la plateforme et présente
les résultats obtenus à chaque étape, leur validation et leur interprétation
financière.

%% ───────────────────────────────────────────────────────────────────────────
\section{Environnement technique}
\label{sec:env}
\par La solution a été développée en \textbf{Python}, avec les bibliothèques
\texttt{pandas} et \texttt{numpy}~\cite{pandas} (manipulation des données),
\texttt{openpyxl} (lecture Excel), \texttt{scikit-learn}~\cite{scikitlearn}
(apprentissage automatique) et \texttt{sqlite3} (entrepôt de données).
L'entrepôt est matérialisé en \textbf{SQLite}, la restitution est prévue sous
\textbf{Power BI}~\cite{powerbi}, et le projet est versionné avec
\textbf{Git/GitHub}. Le code est organisé en modules indépendants orchestrés par
un script unique (\texttt{run\_all.py}), dont l'exécution est illustrée à la
figure~\ref{fig:runall}.

\figplaceholder{Exécution de la chaîne complète (\texttt{run\_all.py}).}{fig:runall}

%% ───────────────────────────────────────────────────────────────────────────
\section{Collecte, nettoyage et unification}
\label{sec:realetl}
\par Les quatre sources ont été unifiées dans un jeu de données long unique. Le
tableau~\ref{tab:volumetrie} en présente la volumétrie. Après suppression des
doublons stricts et correction de \textbf{21 valeurs erronées} (fuites de
colonnes comparatives de l'exercice 2019), le jeu unifié compte \textbf{6\,804
lignes réelles}. La figure~\ref{fig:dataset} en montre un extrait.

\begin{table}[H]
\centering
\caption[Volumétrie des données]{Volumétrie des données par société.}
\label{tab:volumetrie}
\begin{tabular}{|l|c|c|c|c|}
\hline
\textbf{Société} & \textbf{Rôle} & \textbf{Période} & \textbf{Devise} & \textbf{Fréquence} \\ \hline
SFBT & principale & 2005--2025 & TND & semestriel + annuel \\ \hline
DELICE & benchmark & 2015--2025 & TND & annuel \\ \hline
AB~InBev & benchmark & 2015--2025 & USD & annuel \\ \hline
Coca-Cola & benchmark & 2015--2025 & USD & annuel \\ \hline
\end{tabular}
\end{table}

\figplaceholder{Extrait du jeu de données unifié (\texttt{financials\_unified.csv}).}{fig:dataset}

\par La qualité des données a été contrôlée. Le \textbf{contrôle d'équilibre du
bilan} de la \sfbt (Actif = Passif) est vérifié à 0,00~\% sur la grande majorité
des dates, validant la classification automatique des états. Les anomalies
résiduelles (exercices 2017--2019 décalés dans la source, exercices annuels
2021--2022 manquants) sont signalées et documentées plutôt que masquées.

%% ───────────────────────────────────────────────────────────────────────────
\section{Augmentation des données}
\label{sec:realaug}
\par La génération des clôtures de mars et septembre a produit \textbf{3\,320
lignes estimées}, portant le jeu complet à \textbf{10\,124 lignes}. Le
tableau~\ref{tab:augment} illustre la méthode sur le chiffre d'affaires 2024 :
les valeurs estimées (T1, 9M) s'intercalent de façon cohérente entre les valeurs
réelles (S1, FY).

\begin{table}[H]
\centering
\caption[Exemple d'augmentation]{Exemple d'augmentation : chiffre d'affaires SFBT 2024 (M~TND).}
\label{tab:augment}
\begin{tabular}{|l|c|c|c|}
\hline
\textbf{Période} & \textbf{Date} & \textbf{Valeur} & \textbf{Type} \\ \hline
T1 & 31/03/2024 & 183,5 & estimé \\ \hline
S1 & 30/06/2024 & 367,1 & réel \\ \hline
9M & 30/09/2024 & 601,8 & estimé \\ \hline
FY & 31/12/2024 & 836,6 & réel \\ \hline
\end{tabular}
\end{table}

%% ───────────────────────────────────────────────────────────────────────────
\section{Calcul des ratios et du score de performance}
\label{sec:realratios}
\par Les \textbf{46 ratios} du référentiel ont été calculés pour les quatre
sociétés. Le tableau~\ref{tab:ratiossfbt} présente une sélection de ratios clés
de la \sfbt pour l'exercice 2024, comparés à leur fourchette de référence.

\begin{table}[H]
\centering
\caption[Ratios clés de la SFBT]{Sélection de ratios financiers de la SFBT (exercice 2024).}
\label{tab:ratiossfbt}
\begin{tabular}{|l|l|c|c|}
\hline
\textbf{Catégorie} & \textbf{Ratio} & \textbf{SFBT 2024} & \textbf{Benchmark} \\ \hline
Liquidité & Liquidité générale & 3,43 $\times$ & 1,5 à 2,0 \\ \hline
Structure & Autonomie financière & 74,4~\% & 35 à 50~\% \\ \hline
Structure & Endettement global & 25,6~\% & 40 à 60~\% \\ \hline
Structure & Solvabilité & 3,91 $\times$ & 1,5 à 2 \\ \hline
Rentabilité & Marge nette & 30,1~\% & --- \\ \hline
Rentabilité & ROA & 15,8~\% & 5 à 10~\% \\ \hline
Rentabilité & ROE & 26,4~\% & 10 à 15~\% \\ \hline
Risque & Conan \& Holder & 0,21 & $>$ 0,16 \\ \hline
\end{tabular}
\end{table}

\par Ces valeurs dessinent le profil d'une entreprise \textbf{très liquide},
\textbf{peu endettée} et \textbf{fortement rentable}, dépassant la plupart des
fourchettes de référence sectorielles.

\subsection{Score Global de Performance Industrielle}
\par L'agrégation des ratios notés aboutit au \ac{SGPI}. La \sfbt obtient
\textbf{83,5/100} pour l'exercice 2024, devançant nettement ses références
(figure~\ref{fig:sgpi}). Le détail par catégorie révèle une excellence en
liquidité, rentabilité, gestion des risques et productivité, le seul point faible
étant le cycle d'exploitation (stocks élevés).

\begin{figure}[H]
\centering
\begin{tikzpicture}
\begin{axis}[
  width=12cm, height=6cm, ybar, bar width=18pt, ymin=0, ymax=100,
  ylabel={SGPI / 100},
  symbolic x coords={SFBT, AB InBev, Coca-Cola, DELICE},
  xtick=data, nodes near coords, nodes near coords style={font=\footnotesize},
  enlarge x limits=0.18, ymajorgrids, tick label style={font=\footnotesize},
]
\addplot[fill=darkblue,draw=darkblue] coordinates
  {(SFBT,83.5) (AB InBev,55.0) (Coca-Cola,54.5) (DELICE,51.7)};
\end{axis}
\end{tikzpicture}
\caption[Score SGPI par société]{Score Global de Performance Industrielle (SGPI) par société, 2024.}
\label{fig:sgpi}
\end{figure}

\subsection{Analyse comparative (benchmark)}
\par Le tableau~\ref{tab:bench} compare les ratios clés des quatre sociétés. La
\sfbt se distingue par une liquidité et une autonomie financière nettement
supérieures, les références internationales étant plus endettées (effet de
levier). La marge nette de DELICE, supérieure à 100~\%, s'explique par sa nature
de \textbf{holding} (revenus issus de dividendes) et doit être interprétée avec
prudence.

\begin{table}[H]
\centering
\caption[Benchmark des ratios clés]{Benchmark : ratios clés des quatre sociétés (2024).}
\label{tab:bench}
\begin{tabular}{|l|c|c|c|c|}
\hline
\textbf{Ratio} & \textbf{SFBT} & \textbf{DELICE} & \textbf{AB~InBev} & \textbf{Coca-Cola} \\ \hline
Marge nette (\%) & 30,1 & 100,9 & 12,4 & 22,6 \\ \hline
ROE (\%) & 26,4 & 100,0 & 8,4 & 40,4 \\ \hline
Liquidité générale ($\times$) & 3,43 & --- & 0,70 & 1,03 \\ \hline
Endettement global (\%) & 25,6 & 0,4 & 57,1 & 73,8 \\ \hline
\end{tabular}
\end{table}

%% ───────────────────────────────────────────────────────────────────────────
\section{Réalisation de l'entrepôt de données}
\label{sec:realdw}
\par L'entrepôt en étoile a été matérialisé en SQLite. Il comprend quatre
dimensions et des tables de faits (\texttt{Fait\_Etats\_Financiers} : 10\,124
lignes ; \texttt{Fait\_Ratios} : 4\,095 lignes), complétées par des vues
analytiques. L'\textbf{intégrité référentielle} a été vérifiée : aucune ligne de
fait orpheline et aucune perte de donnée lors des jointures. La
figure~\ref{fig:dw} présente l'entrepôt dans un outil d'exploration SQL.

\figplaceholder{Entrepôt de données ouvert dans DB Browser for SQLite.}{fig:dw}

%% ───────────────────────────────────────────────────────────────────────────
\section{Prévision par apprentissage automatique}
\label{sec:realml}
\par Les deux modèles ont été entraînés et évalués par \textit{backtest}
chronologique. Le tableau~\ref{tab:ml} présente leurs performances. La
\textbf{régression linéaire} l'emporte sur l'ensemble des indicateurs, avec un
coefficient de détermination compris entre 0,77 et 0,96 et une erreur (\ac{MAPE})
de 4 à 15~\%.

\begin{table}[H]
\centering
\caption[Performance des modèles de prévision]{Performance des modèles de prévision (backtest, SFBT).}
\label{tab:ml}
\begin{tabular}{|l|l|c|c|}
\hline
\textbf{Indicateur} & \textbf{Modèle} & \textbf{MAPE} & \textbf{$R^2$} \\ \hline
Revenus (CA) & Régression linéaire & 15,3~\% & 0,93 \\ \hline
Total des actifs & Régression linéaire & 4,5~\% & 0,85 \\ \hline
Capitaux propres & Régression linéaire & 4,0~\% & 0,77 \\ \hline
Résultat net & Régression linéaire & 7,9~\% & 0,96 \\ \hline
\end{tabular}
\end{table}

\par Ce résultat est cohérent : la régularité de la croissance de la \sfbt
favorise le modèle linéaire, qui surpasse ici le modèle ensembliste. Un point
méthodologique a été déterminant : modéliser le \textbf{taux de croissance}
plutôt que la valeur brute, ce qui évite l'incapacité des forêts aléatoires à
extrapoler une tendance. La figure~\ref{fig:prevision} confronte les valeurs
réelles aux prévisions.

\figplaceholder{Prévision du chiffre d'affaires : réel vs prédit.}{fig:prevision}

%% ───────────────────────────────────────────────────────────────────────────
\section{Tests et validation}
\label{sec:tests}
\par La fiabilité de la solution est garantie par une \textbf{suite de 70 tests
automatiques} couvrant toutes les étapes : conformité du schéma, unicité des
clés, équilibre comptable, bornes des valeurs estimées, exactitude des ratios
(vérifiés contre des calculs manuels et contre les chiffres publics de Coca-Cola
et AB~InBev), intégrité de l'entrepôt et cohérence des prévisions. La totalité des
tests est satisfaite (\textbf{70/70}), comme l'illustre la
figure~\ref{fig:tests}.

\figplaceholder{Résultat de la suite de tests (\texttt{test\_pipeline.py}).}{fig:tests}

%% ───────────────────────────────────────────────────────────────────────────
\section{Restitution : tableaux de bord}
\label{sec:resti}
\par Les jeux de données produits et l'entrepôt en étoile alimentent des
\textbf{tableaux de bord Power BI} structurés en cinq vues : synthèse \sfbt,
solidité financière, performance, benchmark sectoriel et prévisions. Les
figures~\ref{fig:dashboard1} et~\ref{fig:dashboard2} en présentent un aperçu.

\figplaceholder{Tableau de bord Power BI --- vue d'ensemble de la SFBT.}{fig:dashboard1}

\figplaceholder{Tableau de bord Power BI --- benchmark sectoriel.}{fig:dashboard2}

%% ── Conclusion ─────────────────────────────────────────────────────────────
\section*{Conclusion}
\addcontentsline{toc}{section}{Conclusion}
\par Ce chapitre a présenté la réalisation effective de la plateforme et ses
résultats : des données fiabilisées et enrichies, 46 ratios et un score de
performance plaçant la \sfbt en tête de son secteur, un entrepôt intègre et des
prévisions précises, le tout validé par 70 tests automatiques.
